\documentclass[11pt]{amsart}
\usepackage{mathtools,amssymb,hyperref}
\hypersetup{pdftitle={Capture times, radius bounds and periodic roots in odd-length Kaprekar dynamics},pdfauthor={Rustem Galyavin},hidelinks}
\newtheorem{theorem}{Theorem}
\newtheorem{lemma}[theorem]{Lemma}
\newtheorem{proposition}[theorem]{Proposition}
\newtheorem{corollary}[theorem]{Corollary}
\newtheorem{conjecture}[theorem]{Conjecture}
\newtheorem*{corollary*}{Corollary}
\theoremstyle{definition}
\newtheorem{definition}[theorem]{Definition}
\newtheorem*{remark}{Remark}
\DeclareMathOperator{\sgn}{sgn}
\DeclareMathOperator{\dist}{dist}
\DeclareMathOperator{\tr}{tr}
\DeclareMathOperator{\lcm}{lcm}
\DeclareMathOperator{\ord}{ord}
\newcommand{\Z}{\mathbb Z}
\newcommand{\R}{\mathbb R}
\newcommand{\PP}{\mathbb P}
\newcommand{\one}{\mathbf 1}
\newcommand{\C}{\mathcal C}
\title[Capture times and periodic roots]{Capture times, radius bounds and periodic roots in odd-length Kaprekar dynamics}
\author{Rustem Galyavin}
\date{21 November 2026}
\thanks{\href{mailto:rustem@galyavin.com}{rustem@galyavin.com}. ORCID \href{https://orcid.org/0009-0006-6875-9170}{0009-0006-6875-9170}.}
\subjclass[2020]{11A63, 37E15, 60C05}
\keywords{Kaprekar map, folded doubling, capture time, basin probability, finite automaton}
\raggedbottom
\addtolength{\textheight}{\baselineskip}
\begin{document}
\begin{abstract}
For the Kaprekar map in base \(B\) at a fixed odd length
\(n=2k+1\ge7\), a folded doubling permutation of order \(r\) governs
capture into cycles contained in a chamber defined by digit-pairing
inequalities. At strict lengths the \(r\) gaps on the folded orbit of
slot 1, measured from an anchor gap, evolve in the chamber by an
autonomous integer recurrence whose transient is the time to a cycle.
The probability that a start lies in the chamber, stays there and enters
such a cycle by gap time \(T\le B\) is \(O_n(((T+1)/B)^{k-d})\), where
\(d\) is the dimension of the cone of centred gap profiles that are
constant on the anchor and the folded orbit of slot 1 and satisfy the
secondary chamber comparisons at every folded phase. At strict lengths
\(k-d=r\). At the 47 strict lengths of orders 3 through 13 a matching
lower bound holds for \(0\le T\le c_nB\), with peak radius bounded away
from zero. Every wholly positively paired cycle has normalized peak
radius at most \(2/k+6/B\). For this root recurrence we classify the
loss-free periodic points of period dividing \(2r\) by a fixed finite
graph. An exact full-rank construction extends positive off-centre basin
density to all twelve lengths of folded order 15.
\end{abstract}
\maketitle

\section{Introduction}

The Kaprekar map sorts the digits of a word in descending and ascending
order and subtracts the second number from the first. We keep leading
zeros and sample ordered nonconstant words uniformly. In odd length
\(n=2k+1\), the \(k\) differences of opposite sorted digits determine
the next word. The gaps between these differences have a particularly
simple update where the borrowed digits can be sorted in pairs. The
permutation part of the update is doubling on residues modulo \(n\)
identified with their negatives. Its order is the least \(r>0\) with
\(2^r\equiv\pm1\pmod n\). The corrections to the permutation determine
the time to a cycle.

The mechanism is as follows. Inside the chamber \(\C\), the full-height
gap vectors whose borrowed digits sort in pairs with no secondary
comparison below \(-1\) (Section~\ref{sec:setup}), every correction to
folded doubling falls on the slots
\(S=\{0\}\cup\operatorname{orb}_P(1)\cup\operatorname{orb}_P(3)\), and
the gaps outside \(S\) are only permuted. Call a length \emph{strict} if
it admits a folded-doubling-invariant centred profile, constant on the
root slots \(\{0\}\cup\operatorname{orb}_P(1)\), with every secondary
comparison strictly positive. At a strict length the \(r\) root gaps
measured from the anchor gap \(g_0\) form a relative root
\(x\in\Z^r\). It evolves autonomously by the integer map \(T_r\) of
\cite[Section~2.2]{G9} and \cite[Section~5]{G10}
(equation~\eqref{eq:root}). Three facts make \(r\) the clock.
(i) Let \(K\) be the cone of centred gap profiles that are constant on
the root slots and satisfy the secondary comparisons at all \(r\)
rotated phases, and let \(d=\dim\operatorname{span}K\). States on
chamber cycles lie in a bounded tube around \(B\one/m+K\)
(Lemma~\ref{lem:tube}). After the first folded turn every chamber step
differs from the permutation by a bounded kick, so a start captured
inside the chamber by gap time \(T\) lies at time \(r\) within
\(O_n(T+1)\) of \(B\one/m+\operatorname{span}K\). This tube has
codimension \(k-d\), and bounded numbers of preimages account for the
initial turn, whose kicks can be large. This gives the exponent
\(k-d\) for \emph{chamber capture}, that is, for starts whose orbit up
to and including its first cycle lies in the chamber. At a strict
length \(k-d=r\). Root equality alone gives codimension \(r\), which can miss
constraints. At 25 digits, for example, the rotated secondary comparisons force one further equality,
which reduces the cone to a ray, so \(k-d=11\) while \(r=10\)
(Theorem~\ref{th:clock}). (ii) At a strict
length the gap time from a permanent chamber entrance to the first cycle
equals the transient of \(T_r\) (Theorem~\ref{th:entrance}), and a
finite guard test decides whether a nominal chamber path stays in the
chamber (Theorem~\ref{th:guards}). (iii) Families of \(L^r\) roots whose
transient, spread and loss are \(O(L)\) give the matching lower bound at
the 47 strict lengths of orders 3 through 13
(Theorem~\ref{th:clock}(3)). With \(L\) proportional to \(B\) they give
positive basin mass. We construct such a family at order 15
(Lemma~\ref{lem:cert15}, Theorem~\ref{th:order15}), extending orders
3--13 of \cite{G9}. The clock stops at loss-free periodic points of
\(T_r\). Those of period dividing \(2r\) are classified by one finite
graph (Theorem~\ref{th:graph}), which also gives their one-step
entrances (Proposition~\ref{prop:doors}).

The remaining results mark the chamber's reach. Wholly paired cycles
have normalized peak radius at most \(2/k+6/B\)
(Theorem~\ref{th:radius}). Hence the cut \(\delta_n\) in
Theorem~\ref{th:clock}(3), and the cut produced by the chamber
construction behind Theorem~\ref{th:order15}, are at most \(2/k\).
In the large-base limit the share of starts that begin in \(\C\), leave
it, and exceed a fixed radius cut at some time up to and including their
first state outside \(\C\) decays exponentially in \(n\) when
\(r=o(n)\) (Theorem~\ref{th:entrance}). Starts that stay in \(\C\)
forever have peak at most \(2/k+6/B\), so for large bases none of them
exceeds a fixed cut above \(2/k\). At zero-cone lengths a macroscopic cycle
must lose pairing, or breach a secondary comparison by order \(B\), near
its energy maximum (Theorem~\ref{th:breach}). Digit replication gives
macroscopic fixed points with zero differences, outside the paired
regime, whose directly replicated inputs have negligible original mass
(Proposition~\ref{prop:replicate}). In the large-base limit only a
fraction \((5/8)^{\lfloor(n-1)/4\rfloor}\) of words is positively paired,
and the capture power \(r\) also governs a sector of initially unpaired
starts that enter the chamber in one step (Theorem~\ref{th:sibling}).
Appendix~\ref{app:tails} shows that the chamber exponent does not govern
the full arrival law: at seven digits \(k-d=3\), while the full law has
small-time tail at least \(ct^3\log(1/t)\).
Section~\ref{sec:open} gives the scope of the results and the remaining
population questions.

\subsection*{Related work}

Kay and Downes-Ward use digit multiplicities for fixed points and cycles
in odd and even bases, with modular doubling in the families of
Sections~2 and~5 of the odd-base paper \cite{KDO} and Section~2.2 of the
even-base paper \cite{KDE}. In the odd-base Theorem~7(ii) and
Corollary~8(ii), choosing both scales equal to the replication factor
gives replication on primitive-to-primitive transitions and on primitive
cycles, respectively \cite{KDO}. The proof of their even-base
Theorem~2(a) uses cancellation of the same exceptional borrowed digits
\cite{KDE}. The criterion below gives necessity and sufficiency for
arbitrary digit multisets in any base. Walden determines the seven- and
nine-digit Kaprekar constants, in the sense of fixed points reached from
every nonconstant word, and shows that there are none with 15, 21, 27 or
33 digits \cite{Walden}. Chen, Ono, Schwartz and Thakur give a conjugacy
to projective doubling modulo the base on a forward-invariant region for
four digits in odd bases greater than three \cite[Theorem~1]{COST}.
Dahl studies convergence and projected gap
dynamics at fixed decimal lengths \cite[Sections~2--3]{Dahl}.
Devlin and Zeng give four-digit basin and maximum-distance results
\cite[Theorems~1--2]{DZ}. Further cycle and arrival results in even
bases are in \cite{G1,G2}. Thakur records the difference coordinates and borrowed digits, and
discusses base translations for cycles \cite[Section~2.7 and Remark~4.3]{Thakur}.
We use the full limiting laws at five and seven digits \cite{G3,G5}, and the gap majorization and exact base
covariance of \cite{G4}. At eleven digits, \cite{G6,G7,G8} treat
access, arithmetic selection and an entrant law.
The paired chamber update, strict profiles and full-rank source
constructions come from \cite{G9,G10}. The latter gives the chamber
criterion and the full-orbit case of the breach bound, while \cite{G11}
studies the arithmetic of the invariant cone.

The sign feedback in the recurrence is analogous to the use of quantized
feedback in sigma--delta modulation \cite[Section~3.2, equation~(31)]{DD}.
For the one-bit modulators they study, Reefman, Reiss, Janssen and
Sandler test periodic output words by state-space closure and quantizer
inequalities. For the architectures with a dc pole, their balance
condition equates average input and output \cite[Sections~II--III,
equations~(5), (9) and (15), and Appendix~A]{RRJS}. The same authors
extend this description to feedback modulators \cite{RRJS2}, and Perez
Gonzalez and Reiss separate dc idle tones from long and short limit
cycles in simulations \cite{PGR}. Akiyama and Peth\H{o} show that a
discretized rotation of the integer lattice has infinitely many periodic
orbits \cite{AP}. We use none of these results; the comparison is an
analogy. Here the feedback is ternary, the recurrence has an integer
loss term, and its period-\(2r\) closure gives opposite signs across the
two turns. That specific identity yields the finite graph and the count
proved below.

\emph{Contributions.} From \cite{G9,G10} we use the paired update
\cite[Lemma~4]{G10}, the root map \(T_r\) \cite[Section~2.2]{G9},
\cite[Section~5]{G10}, the retirement construction
\cite[Section~4]{G9}, the chamber criterion \cite[Theorem~1]{G10},
strict source transport \cite[Proposition~7]{G10}, and the strict family
at \(2^a\pm1\) and the composite vector \cite[Section~3]{G10}; from
\cite{G4} we use majorization and the paired covariance. New here, as
far as we know for Kaprekar dynamics, are the cone \(K\) with the
capture exponent \(k-d\), the breach bound at zero-cone lengths that are
not full-orbit, the paired radius ceiling, the replication criterion
(special cases of which are in \cite{KDO,KDE}), the initial pairing
limits and reflected entrances, the entrance comparison and the
exceedance bound for initial chamber segments, the read recurrence with phase balance and the guard test, the
periodic root graph with its corollaries, and the order-15 records.

\section{Gap coordinates and original measure}\label{sec:setup}

For descending sorted digits \(a_1\ge\cdots\ge a_n\) put
\[
 \begin{aligned}
 q_i&=a_i-a_{n+1-i}\quad(1\le i\le k),\\
 g&=(B-q_1,q_1-q_2,\ldots,q_{k-1}-q_k,q_k).
 \end{aligned}
\]
Thus \(g_0\ge1\), \(g_j\ge0\) for \(j\ge1\), and \(\sum_jg_j=B\).
These conditions characterize the physical gap vectors. For example,
one may realize the differences by putting all the smaller digits and
the middle digit at zero. We write \(g_t\) for the vector at time \(t\)
and \((g_t)_j\) for a coordinate. Let
\[
 R(g)=\max_{1\le i\le k}|q_i/B-(m-i)/m|,\qquad m=k+1,
\]
and let \(Z_B\) be the maximum of \(R\) over all phases of the first cycle.
A positive threshold independent of \(B\) marks a macroscopic radius.
The next numerical word is \(\sum_{i=1}^k q_i(B^{n-i}-B^{i-1})\).
Consequently the gap and word cycles have the same least period. The
first cyclic word occurs at the first cyclic gap time or one step later.

Let \(\theta\) be the number of positive differences, called the height.
For \(\theta>0\), column subtraction gives a provisional multiset of
\(n-2\theta\) copies of \(B-1\) and the pairs \(q_i,B-1-q_i\) for
\(i\le\theta\).
The exact multiset decreases \(q_\theta\) by one and increases the reflected
digit \(B-1-q_1\) by one. These are the innermost direct column reached
by the borrow and the rightmost column, respectively. At height zero
the output is the all-zero word.

From now on \(n\ge7\) is odd unless specified otherwise. At full
height \(\theta=k\) the singleton \(B-1\) is followed by two columns
\[
 A_i=q_i-\one_{i=k},\qquad C_i=B-q_{m-i}-\one_{i<k}\quad(1\le i\le k).
\]
Positive pairing means \(q_k>0\) and
\(\min(A_i,C_i)\ge\max(A_{i+1},C_{i+1})\) for every \(i<k\).
Write
\[
 \begin{gathered}
 p_0=\lfloor m/2\rfloor,\qquad 1\le i\le p_0,\\
 d_i(g)=\sum_{j=m-i}^{m-1}g_j-\sum_{j=0}^{i-1}g_j,\qquad
 h_i=(-1-d_i(g))_+ .
 \end{gathered}
\]
We use \(d_0=h_0=0\). The chamber \(\C\) consists of positively
paired gaps with \(h_i=0\) for \(i\ge2\).
The permutation \(P\) sends slot \(j>0\) to \(\min(2j,n-2j)\) and
fixes 0. It has order \(r\) because the orbit of 1 has that order
and \(2^r\equiv\pm1\pmod n\) also fixes every other slot. Write
\(J=\{0\}\cup\operatorname{orb}_P(1)\). We have \(P(k)=1\) and
\(P(k-1)=3\).

\begin{lemma}[Paired update {\cite[Lemma~4]{G10}}]\label{lem:update}
On a positively paired source,
\[
 g^+=Pg+a+ev,\qquad e=\sgn(d_1(g)),\qquad v=(1,-2,1,0,\ldots),
\]
where \(a\) has successive pairs \((-\Delta h_i,\Delta h_i)\),
\(\Delta h_i=h_i-h_{i-1}\), and an unmatched last entry zero if \(m\) is odd.
In \(\C\) this gives
\begin{equation}\label{eq:Cupdate}
 g^+=Pg+ev+h_1(-1,1,1,-1,0,\ldots).
\end{equation}
\end{lemma}

\begin{proof}
We recall the argument of \cite[Section~2]{G10}.
For the provisional calculation replace the last pair by
\((q_k,B-1-q_1)\). Keep the given level order, taking the larger and
smaller entries within each pair.
The sum of the two resulting gaps in slots \(2i-2,2i-1\)
is \(g_{i-1}+g_{m-i}\), and their difference (first minus second)
is \(|1+d_{i-1}|-|1+d_i|\). The unmatched gap is \(g_{k/2}\). Since
\(|1+d_i|=1+d_i+2h_i\), the resulting gaps are \(Pg+a\).
Restoring the last level changes its maximum and minimum by
\(-e,e\), respectively, since on integers
\(\min(g_k,g_0+1)-\min(g_0,g_k+1)=e\).
This changes the first two differences by \(-e,e\) and the gap
vector by \(ev\).
\end{proof}

Let \(\PP_{n,B}\) denote original probability on the \(B^n-B\)
ordered nonconstant words. The distinct-digit multiplicity of a positive
gap vector is \(n!g_0\prod_{j=1}^k(g_j-1)\).
The smallest digit has \(g_0\) choices and each pair of opposite
sorted spacings has \(g_j-1\) positive splits. Weak splits give the upper
bound \(n!g_0\prod_{j=1}^k(g_j+1)\) on all realizations.
Thus the weight of a state is \(O_n(B^{-k})\), and it is
\(\asymp_n B^{-k}\) uniformly on a fixed compact subset of the
open normalized simplex. Coupling the original digits to independent
uniforms on \([0,1]\) by flooring \(BU\) shows that the normalized gaps
converge to \(\mathrm{Dirichlet}(2,\ldots,2)\). Indeed the \(n+1=2m\)
real spacings, including both end spacings, are uniform on their simplex.
Summing spacings in opposite pairs gives the density
\(n!\prod_{j=0}^ku_j\) of the normalized gap vector \(u\). Conditioning
on nonconstancy has asymptotically no effect.

\emph{Terms.} Slot 0 is the \emph{anchor}, and the gaps in \(J\) are
the \emph{root gaps}. The comparison \(d_1\) is the \emph{primary} row
and the \(d_i\) with \(i\ge2\) are the \emph{secondary} rows; \(h_i\) is
the \emph{breach} of row \(i\). Thus \(\C\) consists of the positively
paired gaps with no secondary breach. A length is \emph{strict} if some
folded-doubling-invariant centred profile, constant on \(J\), has every
secondary row strictly positive (Section~\ref{sec:clock}). A
\emph{chamber cycle} is a cycle all of whose phases lie in \(\C\). A
\emph{permanent chamber entrance} of an orbit is a time \(t_e\) with
\(g_t\in\C\) for all \(t\ge t_e\). The \emph{nominal} path of a chamber
state iterates \eqref{eq:Cupdate} without checking agreement with
subtraction. Its \emph{guards} are the inequalities for physical gaps,
full height, positive pairing and \(h_i=0\) for \(i\ge2\).
\emph{Marked} capture is capture into a cycle with \(Z_B\) above a
stated cut. The \emph{lifetime norm} of a path \(x(t)\) in \(\Z^r\) is
\(\sup_t\|x(t)\|_\infty\).

\section{The capture exponent and the breach bound}\label{sec:clock}

In the zero-sum gap space define the rational cones
\[
 K=\{u:u_j=u_0\ (j\in J),\quad
 d_i(P^t u)\ge0\ (2\le i\le p_0,\ 0\le t<r)\}.
\]
Let \(W=K\cap\operatorname{Fix}(P)\). Set \(d=\dim\operatorname{span}K\) and
\(d_f=\dim\operatorname{span}W\). A length is \emph{strict} if
\(W\) contains a direction with every secondary row strictly positive.
Equivalently there is a normalized invariant profile \(\pi\), constant
on \(J\), with those positive slacks. By mixing with the uniform
profile one can take \(\pi\) sufficiently close to it that all homogeneous
pairing comparisons have positive margin. Indeed, for uniform gaps the
two unborrowed columns coincide at every level and successive levels
differ by \(1/m\).

Let \(A_B(T)\) consist of starts in \(\C\) which stay there and enter
a cycle contained in \(\C\) by gap time \(T\).
Thus \(A_B(0)\) requires the initial gap to lie on such a cycle.

\begin{theorem}[Dimension and capture time]\label{th:clock}
At fixed odd length \(n\ge7\), the following hold.
\begin{enumerate}
\item For all integers \(0\le T\le B\),
\(\PP_{n,B}(A_B(T))\ll_n ((T+1)/B)^{k-d}\).
The same bound holds, with a constant depending also on \(j\), for
starts whose gaps from time \(j\) on stay in \(\C\) and enter a cycle
contained in \(\C\) by gap time \(j+T\).
\item In each class modulo \(m\) containing a base with a chamber cycle,
\(\PP_{n,B}(A_B(0))\asymp_n B^{d-k}\) for large bases in the class.
At strict lengths \(k-d=r\). In particular at \(n=2^a\pm1\), \(a\ge4\),
the exponent is \(a\).
\item At each of the 47 strict lengths of folded orders 3 through 13
in \cite{G9}, there exist \(\delta_n,c_n>0\) such that for all
large integer bases and \(0\le T\le c_n B\),
\[
 \PP_{n,B}(A_B(T),Z_B>\delta_n)\asymp_n ((T+1)/B)^r.
\]
\item Fix a numerical-word time \(T\ge1\). With arbitrary prefixes,
the probability of reaching a chamber cycle by that time is
\(O_{n,T}(B^{d-k})\).
For fixed points the corresponding exponent is \(d_f-k\).
If \(d>0\), for sufficiently small \(\delta>0\) the probability of
reaching by that time a chamber cycle of peak \(Z_B>\delta\) has the
matching lower order on the base multiples of the positive
generator of \(\{\sum_j a_j:a\in
(\R\one+\operatorname{span}K)\cap\Z^m\}\).
For fixed points replace \(K,d\) by \(W,d_f\).
The generator is 1 at strict or composite lengths. At 17 digits the
respective probability exponents are \(-4,-7\).
\item At 25 digits \(K\) is a ray and \(k-d=11\), while \(r=10\).
\end{enumerate}
\end{theorem}

\begin{remark}
The exponent concerns chamber capture only. Seven is a full-orbit length
(\(J\) contains every slot), so \(K=\{0\}\) and \(k-d=3\). Thus (1)
gives \(\PP_{7,B}(A_B(\lfloor tB\rfloor))\ll t^3\) for
\(B^{-1}\le t\le1\), while the full arrival law satisfies
\(F_{7,\beta}(t)\ge ct^3\log(1/t)\) (Proposition~\ref{prop:tails}). For
small \(t\), in the large-base limit, most orbits captured by time
\(tB\) are therefore not chamber-resident from the start. For arbitrary
prefixes we prove only the fixed-time statement (4), the statement in
(1) after a fixed number of unconstrained steps, and the one-step sector
of Theorem~\ref{th:sibling}.
\end{remark}

\begin{theorem}[Breach at a zero invariant cone]\label{th:breach}
Fix an odd length with \(W=\{0\}\), and choose a phase of a physical
cycle maximizing \(V(g)=\|g-B\one/m\|_2^2\). In the preceding
\(2r\) source phases, counted cyclically, either some source fails
positive pairing or
\[
 BZ_B\le C_n(1+H),
 \qquad H=\max_{\text{those phases},\ 2\le i\le p_0} h_i .
\]
Here \(C_n\) depends only on \(n\).
In particular a cycle of radius at least \(\delta>0\) that pairs
through the window requires \(H\ge\delta B/C_n-1\).
\end{theorem}

At full-orbit lengths, where \(J\) contains every slot, this is
\cite[Theorem~2]{G10}. The new cases are the lengths with \(W=\{0\}\)
that are not full-orbit, such as 229, 283 and 1069
(Section~\ref{sec:verify}). The energy estimate is that of \cite{G10}.
The cone-residual bound and orbit averaging carry its control from the
root slots to all slots.

\begin{proof}[Energy estimate and proof of Theorem~\ref{th:breach}]
In any \(2r\) consecutive positively paired steps with secondary
breaches bounded by \(H\), let \(g\) be the midpoint and set
\(u=g-B\one/m\). Write \(z=g_k-g_0\) at an arbitrary source.
The next slot 1, relative to the anchor, is
\(z-3\sgn z+2(-1-z)_+\ge-2\).
An update of any other relative slot decreases its permuted value
by at most \(1+2H\), since the primary correction is nonnegative
there and the secondary correction vanishes at the anchor.
The slot feeding \(k\) passes through slot 1 within one turn.
As in the proof of \cite[Theorem~2]{G10}, following slot \(k\in J\)
over the last \(r\) updates gives \(h_1\le1+(r-1)(1+2H)\) there.
Hence on the last turn the primary correction and the total defect
from \(P\) are \(O_n(1+H)\). If the last turn starts at \(g\), at
its phase \(t\) the state differs from \(P^t g\) by \(O_n(1+H)\).

By the paired update the sorting work is
\[
 2\langle Pg,a\rangle+\|a\|_2^2
   =2\sum_i\Delta h_i\,\Delta(d_i+h_i)\le0,
\]
since \(h_i=(-1-d_i)_+\) is nonincreasing in \(d_i\) while
\(d_i+h_i\) is nondecreasing. The signed borrowing work against
\(Pg\) is \(-4|g_k-g_0|+2e(g_1-g_0)\).
Along the ideal last-turn permutation put \(b_t=(P^tg)_k-g_0\).
Because \(P(k)=1\), the signed work is at most
\(-4|b_t|+2|b_{t-1}|+O_n(1+H)\), with indices taken cyclically.
The sorting-borrowing cross term and borrowing square are also
\(O_n(1+H)\) per step. Summing gives
\begin{equation}\label{eq:work}
 V_{\rm end}-V_{\rm mid}\le -2\sum_{j\in J\setminus\{0\}}
 |u_j-u_0|+O_n(1+H).
\end{equation}
The secondary violations \([-d_i(P^tu)]_+\) are likewise
\(O_n(1+H)\) by the last-turn comparison.

For the cones used here, distance to the cone is at most a fixed
constant times the sum of absolute equality residuals and negative
parts of the required inequalities. This is Hoffman's error bound for
linear inequalities \cite{Hoffman52}; we include the short argument for
these cones. Subdivide the
zero-sum space by the signs of those linear forms and the coordinates.
On each cone the residual sum is a nonnegative linear form. The
vertices of its intersection with the unit \(\ell^1\) sphere
generate the cone. In a nonnegative expansion discard the generators
having positive residual. The retained sum satisfies all constraints,
and the discarded norm is at most a constant times the residual sum.
The constant is finite because the collection of generators is finite.
This argument also applies after including the equalities \(Pu=u\).

Averaging the orbit of a point of \(K\) gives a point of \(W\).
If the average is zero then every rotated secondary row vanishes
because all are nonnegative. The primary row vanishes by root equality.
The zero rows give reflected pairs of equal gaps at each rotated phase.
On such vectors the value at slot \(j\) after applying \(P\)
is the previous value at \(\lfloor j/2\rfloor\). After
\(\lceil\log_2m\rceil\) phases it is constant, so invertibility
and zero sum give \(u=0\). Thus \(K=\{0\}\) if and only if
\(W=\{0\}\). This is the orbit-averaging argument in the proof of
\cite[Theorem~1]{G10}.

If the window ends at a global energy maximum, its last-turn work
is nonnegative. Equation~\eqref{eq:work} and the secondary residual
bounds give \(\dist(u,K)=O_n(1+H)\). The endpoint differs
from the midpoint by \(O_n(1+H)\) because \(P^r=I\).
At \(K=\{0\}\) the endpoint has centred norm \(O_n(1+H)\),
and \(BZ_B\le \sqrt{mV_{\rm end}}\) by the suffix sums and maximality.
\end{proof}

\begin{lemma}[Cyclic tube and bounded preimages]\label{lem:tube}
Every state on a cycle in \(\C\) lies within \(O_n(1)\) of
\(B\one/m+K\). For fixed points the same holds with \(W\).
Every target gap vector, including those reached from sources of
height below \(k\), has at most
\(1+\sum_{\theta=1}^k3^\theta(2\theta)!\) gap predecessors.
\end{lemma}

\begin{proof}
The first claim is the root-diameter step in the proof of
\cite[Theorem~1]{G10}, followed by the cone residual bound.
In the chamber, if slot 3 is outside \(J\), the complementary sum
decreases by \(h_1\) each step. On a cycle this forces \(h_1=0\).
If 3 belongs to \(J\), all corrections already act there. Thus
on a chamber cycle the complementary gaps permute and the sum on \(J\)
is preserved. Hence along the cycle \(V\) and the energy on \(J\) about
its own mean differ by a constant, and they have the same maximizing
phases. When \(3\in J\), the primary correction is bounded on the last
turn by the slot-1 reset over the preceding turn.
The last-turn work calculation with \(H=0\) applies
to the energy on \(J\) around its own mean, with the same terms because
the corrections are supported there. Choose a maximum of this energy
as the endpoint. Nonnegative work bounds the midpoint anchor
differences on \(J\), and the endpoint is within \(O_n(1)\) of the
midpoint. Maximality then gives a bounded root diameter at every
phase. All kicks on the cycle are bounded, and every phase is within
one bounded turn of its ideal permutation. The cone residual bound
gives the first claim. A fixed point also has \(Pu-u=O_n(1)\),
giving the claim for \(W\).

For the inverse count fix a positive height \(\theta\) and let
\(a=n-2\theta\ge1\). The \(2\theta\)-digit borrowed core has a matching
into pairs whose sums differ from \(B-1\) by at most one.
This also holds when both borrow corrections affect one pair.
Lowering one element of a multiset by one lowers exactly one entry of
its sorted list by one, and raising one element raises exactly one
entry. So the sorted core is a reflection-symmetric list with one entry
lowered and one raised. In the descending output list \(z\), fix its
first \(a\) entries to \(B-1\). Hence
\(z_{a+i}+z_{n+1-i}=B-1+\varepsilon_i\) with
\(\varepsilon_i\in\{-1,0,1\}\); guess these \(\theta\) mirror errors.
From the target differences derive successively
\[
 z_{n+1-i}=z_i-q_i',\qquad
 z_{a+i}=B-1+\varepsilon_i-z_{n+1-i}\quad(1\le i\le\theta).
\]
The needed \(z_i\) is already known since \(a\ge1\).
Placing the core digits in the borrowed positions uses at most
\((2\theta)!\) permutations. This recovers the positive differences
from their direct digits. Height zero adds at most one predecessor.
\end{proof}

\begin{proof}[Proof of Theorem~\ref{th:clock}, upper bounds and static lower bounds]
After the first \(r\) chamber sources all defects from \(P\)
are \(O_n(1)\) by the slot-1 reset in the energy proof. If the
path first enters its cycle after time \(r\), reverse the
permutation and sum the defects from that entry back to time \(r\).
The space \(\operatorname{span}K\) is \(P\)-invariant, so the
time-\(r\) state is within \(O_n(T+1)\) of
\(B\one/m+\operatorname{span}K\). If entry is earlier, the
state at time \(r\) is already in the cyclic tube. Bounded
fundamental cells in the sum-\(B\) lattice give a point count
\(O_n(B^d(T+1)^{k-d})\) by thickening in parallel and normal
directions. The inverse count through the first \(r\) steps, and
the original gap weights, prove the first bound in (1). Applied from
time \(j\), with the inverse count through the first \(j+r\) steps, the
same argument gives the prefix statement in (1), with a constant
depending also on \(j\).

For (2), take a seed \((g_s)\) at base \(b\). At base
\(B=b+mL\), perturb phase \(s\) by \(L\one+P^s w\), with
integral \(w\in K\) and \(\|w\|_\infty\le\epsilon L\).
The uniform translation increases every interlevel margin by \(L\),
as in the exact covariance of \cite[Section~6]{G4}.
The remaining perturbation, of size \(O_n(\epsilon L)\), leaves the
primary row unchanged and increases or preserves every secondary
row. Thus the paired correction remains unchanged and the translated
phases close after a common multiple of the seed period and \(r\).
A fixed relative box in the cone dilates to \(\gg_n L^d\)
integral points, each with gap weight \(\asymp_n B^{-k}\).
If \(d>0\) the box can avoid zero so the cycles have a fixed
positive peak bound, since suffix sums give a norm on the zero-sum
space. The same cutoff can be used over the finitely many seeded
classes. If \(d=0\) the uniform translation alone suffices.
At strict lengths the open cone relative to the root-equal plane
has dimension \(k-r\). Such lengths have seeds in every class:
round a near-uniform strict invariant profile with common root
rounding and adjust the sum at a complementary slot, keeping the
slacks in all rotations for large \(B\). For \(2^a+1\) the orbit of 1 consists of the powers of two
through \(2^{a-1}\). For \(2^a-1\) the last slot is
\(2^{a-1}-1\) in place of \(2^{a-1}\). Both orbits have
length \(a\). Assigning weight zero on \(J\) and one elsewhere
makes each secondary row positive. The prefix contains at
least two root-or-anchor slots and the suffix only one, except
at the longest minus-case row where the prefix contains at
least three and the suffix contains two. This is the strict family of \cite[Section~3]{G10}.

For (4), at time \(T\) the gap is already cyclic. The tube count
is \(O_n(B^d)\), or \(O_n(B^{d_f})\) for fixed points. The number
of preimages under \(T\) steps is bounded uniformly in \(B\) by
Lemma~\ref{lem:tube}. For the lower bound let
\(S_0=\R\one+\operatorname{span}K\) and let \(b\) be the generator
in the statement. At each multiple of \(b\), integral points of
\(S_0\) in the sum-\(B\) plane form a translate of the full
zero-sum lattice in \(\operatorname{span}K\).
A dilated box in a small relative cone slice away from zero gives
\(\gg_n B^d\) such gaps near the uniform vector. Their rotations
pair, the corrections vanish, and each original at such a gap
maps to a cyclic word in one step. The fixed-point proof uses \(W\).

At a strict length \(S_0\) contains a unit vector in a
complementary slot. For the fixed subspace the sums include
\(r+1\) from \(J\) and the sizes of the complementary orbits,
which divide \(r\). For an odd composite and a proper divisor
\(1<\gamma<n\), the vector \(v_\gamma(j)=\one_{j>0,\ \gamma\mid j}\) is
invariant and zero on \(J\) (for prime \(\gamma\) this is the vector of
\cite[Section~3]{G10}). The images under \(P\) of the last \(i\) slots
are \(1,3,\ldots,2i-1\). Counting multiples of \(\gamma\) gives
\[
 d_i(v_\gamma)=\one_{(i-1)\bmod\gamma\,\ge(\gamma-1)/2}.
\]
So \(v_\gamma-\overline v_\gamma\one\in W\). The integer sums for the
fixed span with constants include \(m\) and, using \(\gamma=n/\ell\),
all \((\ell-1)/2\) for primes \(\ell\mid n\). Their gcd is 1:
any prime dividing all the latter gives \(n\equiv1\pmod{2q}\)
for that prime \(q\), and hence cannot divide \(m\).
The same conclusion holds for the larger space.
At length 17 the complement is one orbit with four slots,
so \(d=4,d_f=1,k=8\).

At length 25 only slots 5 and 10 lie outside \(J\).
Let their offsets from the root be \(x,y\). The secondary rows
are \(0,y,y,y,y-x\), and a turn exchanges the two offsets at its
first step. Thus \(K\) is the ray \(x=y\ge0\) with zero sum.
Take every gap equal to \((B-2t)/13\) except for an additional
\(t\) at 5 and 10, with \(t\) near \(B/400\) and
\(2t\equiv B\pmod{13}\). This is a fixed gap with positive
macroscopic radius at every sufficiently large base.
\end{proof}

\subsection{Marked lower bound}\label{sec:transport}

At a strict length slot \(k-1\) is outside \(J\):
on an invariant root-equal profile \(d_2=g_{k-1}-g_0\),
and it must be positive. Thus slot 3 is also outside \(J\).
List the nonzero root slots in orbit order as \(j_1,\ldots,j_r\),
with \(j_r=k\), and put \(x_i=g_{j_i}-g_0\).
In the chamber, \eqref{eq:Cupdate} gives the autonomous root map of
\cite[Section~2.2]{G9} (see also \cite[Section~5]{G10}):
\begin{equation}\label{eq:root}
 \begin{gathered}
 T_r(x)=(z-3e+2h,\ x_1+2h,\ x_2-e+h,\ldots,x_{r-1}-e+h),\\
 z=x_r,\qquad e=\sgn z,\qquad h=(-1-z)_+ .
 \end{gathered}
\end{equation}
We call \(x\) the relative root and \(h\) its loss. The absolute
sum on \(J\) increases by \(h\), while the complement permutes with
a subtraction of \(h\) in slot 3. A root path is loss-free
when \(h=0\) at every step.

The finite records of \cite[Appendix~B]{G9} give the following
input to the proof. Appendix~\ref{app:retirement} explains how
a finite record gives the uniform family.
\begin{lemma}[Full-rank source records of orders 3--13]\label{lem:p9}
For each \(3\le r\le13\) and integer \(L\ge1\), the supplied
integer records for \eqref{eq:root} give \(L^r\) distinct starts.
Their time through a loss-free terminal return, lifetime spread
\[
 \sup_t\bigl[\max(0,x_1(t),\ldots,x_r(t))-\min(0,x_1(t),\ldots,x_r(t))\bigr]
\]
and total loss are bounded respectively
by \(C_{r,\mathrm{time}}L,C_{r,\mathrm{sp}}L,C_{r,\mathrm{loss}}L\)
for constants independent of \(L\).
The associated length records give strict profiles for all 47
strict lengths in that order range.
\end{lemma}

\begin{proof}[Proof of Theorem~\ref{th:clock}(3)]
This is the strict source transport of \cite[Proposition~7]{G10} with
the time budget of Lemma~\ref{lem:p9} tracked.
Fix a near-uniform strict invariant profile \(\pi\) and take
\(L=\lfloor T/(2C_{r,\mathrm{time}})\rfloor\) for
\(T\ge4C_{r,\mathrm{time}}\).
For each relative start \(x\), choose the anchor in an interval
of \(\lfloor\epsilon B\rfloor\) integers starting at
\(\lfloor B\pi_0-\sum_i x_i/(r+1)\rfloor\).
Choose all but one complementary gap independently over intervals
of the same size starting at their profile floors.
The last gap enforces the sum \(B\), at every integer base.
This gives \(\gg_n B^{k-r}\) choices for each root start.
Along the nominal chamber trajectory the mean on \(J\) changes
by total loss divided by \(r+1\), and all complementary errors
from pure permutation are bounded by total loss. Together with
the spread budget this keeps the trajectory within
\(O_n(1+\epsilon B+T)\) of \(B\pi\).
Small enough \(\epsilon,c_n>0\) preserve positive gaps, pairing,
secondary slacks and a positive radius cut when \(T\le c_n B\).
Thus the nominal trajectory is physical. At root cycle entry the
root sum is constant thereafter, and the complement permutes
without losses. The absolute gaps are recurrent from that
entry, with a return at a common multiple of the root period
and \(r\).

The interior starts have total original weight
\(\gg_n (L/B)^r\).
For bounded smaller \(T\), the static marked construction works
in every class and gives the same order, after reducing the
common cut and constant. The upper bound is (1).
\end{proof}

\section{Radius bounds on paired paths}\label{sec:radius}

\begin{theorem}[Paired radius and original prefix]\label{th:radius}
If the next \(N\) source states starting at \(g\) are positively
paired, then \(R(g)\le2^{-N}+2/k+6/B\).
Thus every wholly positively paired cycle satisfies
\(Z_B\le2/k+6/B\), even with secondary breaches.
Along lengths \(n\equiv3\pmod6\) the order \(1/k\) is sharp up to
constants.
For each \(0<\epsilon<1/2\) and \(\delta>0\),
\[
 \lim_{\substack{n\to\infty\\n\text{ odd}}}\limsup_{B\to\infty}
 \PP_{n,B}\!\left(\max_{0\le t\le\lfloor(1/2-\epsilon)\log_3 n\rfloor}
 R(g_t)>\delta\right)=0.
\]
\end{theorem}

\begin{proof}
For \(\mu=k^{-1}\sum_i\delta_{q_i/B}\) let \(D(\mu)\) be the
uniform CDF discrepancy. The target \(j/(k+1)\) of the \(j\)-th
increasing atom lies between its empirical left and right
ranks, so \(R\le D\le R+1/k\). At a tie the extreme ranks of
the tied block give the two CDF bounds used here.

An allowance \((u,v)\) between two measures means that each CDF
at \(x\) is at most the other at \(x+u\) plus \(v\).
Allowances compose by addition and bound the discrepancy
difference by \(u+v\). Let \(s(x)=1-x\), \(L(x)=|2x-1|\), and
\(Q\nu=\tfrac12((S_-)_\#\nu+(S_+)_\#\nu)\) with
\(S_\pm(y)=(1\pm y)/2\), where \(\#\) denotes pushforward.
Then \(QL_\#\mu=(\mu+s_\#\mu)/2\) and \(D(Q\nu)\le D(\nu)/2\).
For the latter the CDF on either half is half an affine
function of the original CDF at \(1-2x\) or \(2x-1\),
using a left limit on the decreasing branch.
Define \(b(q)\) as the least vertical allowance comparing
\(\mu\) with its reflection at horizontal allowance \(2/B\).
The allowance between \(\mu\) and \(QL_\#\mu\) is therefore
\((2/B,b(q)/2)\).

To compare with the actual successor, sort \(L(q_i/B)\) in
descending order as \(d_1',\ldots,d_k'\), and put \(d'_0=1\).
The symmetric differences of the sorted ideal digits
\(1,q_i/B,1-q_i/B\) are \((d'_{i-1}+d'_i)/2\).
They interlace the \(d'_i\), giving vertical allowance \(1/k\).
If \(z\) differences vanish, form the true borrowed digits from
the ideal list by replacing \(z\) low zero-pair members with
\(1-1/B\) and changing the other matched digits by at most
\(1/B\). In descending order every new rank \(j\) is between
the old ranks \(j+z,j-z\) up to \(1/B\), using just the side
available when a rank falls outside the list.
Subtracting the resulting outer and inner bounds puts the new
difference at rank \(i\) between the old differences at
\(i+z,i-z\) up to \(2/B\). Thus \(L_\#\mu\) and the successor
\(\mu^+\) have allowance \((2/B,(z+1)/k)\).
At height zero the same bound holds because its vertical part
exceeds one. Consequently
\[
 D(\mu)\le\tfrac12 D(\mu^+)
       +\tfrac12\big(b(q)+z/k+1/k\big)+3/B.
\]
For a positively paired source each threshold separates the
two borrowed column counts by at most one, since only one level
can straddle the threshold. Each column is within \(1/B\) of
its ideal list, so \(b(q)\le1/k\) and \(z=0\).
Iteration proves the radius bound. For a general path it yields
\[
 R(g_t)\le2^{-N}+\max_{t\le j<t+N}\{b(q(j))+z_j/k\}+1/k+6/B.
\]
If \(R(g_t)>\delta\) and the sum of the three explicit error
terms is at most \(\delta/2\), the window contains a source
with \(b(q)>\delta/4\) or \(z/k>\delta/4\).

Finally \(D(L_\#\mu)\le2D(\mu)\), since the preimage of a
CDF interval is a closed interval of corresponding uniform
length. The allowance and \(z/k\le D(\mu)\) give
\(D(\mu^+)\le3D(\mu)+1/k+2/B\).
If the original uniforms have empirical discrepancy at most
\(a\), their order statistics are within \(a+1/n\) of
\(j/(n+1)\). Rounding gives \(D(\mu_0)\le2a+2/n+1/k+1/B\).
By the Dvoretzky--Kiefer--Wolfowitz inequality with Massart's
constant \cite[p.~1269]{Massart}, the exceptional probability
is at most \(2e^{-2na^2}/(1-B^{1-n})\) under original sampling.
With \(a=n^{-1/2+\epsilon/2}\), iteration up to the stated horizon
proves the iterated limit.
\end{proof}

The window estimate in the proof gives the following quantitative form
for arbitrary paths.

\begin{corollary*}[Window form]
For a difference vector \(q\) let \(\mu_q=k^{-1}\sum_i\delta_{q_i/B}\),
and let \(b(q)\) be the least \(v\ge0\) such that, for every \(x\), the
CDF of each of \(\mu_q\) and its reflection under \(x\mapsto1-x\) at
\(x\) is at most the other at \(x+2/B\) plus \(v\). Let \(\delta>0\) and
\(N\ge1\) satisfy \(2^{-N}+1/k+6/B\le\delta/2\). If \(R(g_t)>\delta\) on
an arbitrary path, then some \(j\) with \(t\le j<t+N\) has
\(b(q(j))>\delta/4\) or more than \(\delta k/4\) zero differences.
\end{corollary*}

\begin{proof}
This is the window estimate at the end of the third paragraph of the
preceding proof, where \(z_j\) is the number of zero differences at
source \(j\).
\end{proof}

For sharpness take \(n=6v+3\ge9\) and
\(g_j=2L+L\one_{j>0,\ 3\mid j}\), at \(B=(2m+v)L\).
This invariant profile is constant on \(J\), and counting
multiples of 3 as in Section~\ref{sec:clock} gives
\(d_i/L\in\{0,1\}\) and \(d_1=0\).
The unborrowed digits within each level differ by \(d_i\)
(using \(d_{m-i}=d_i\)), and each column drops by at least
\(2L\) between levels. The pairing margins after borrowing
are at least \(L-1\). All corrections in \eqref{eq:Cupdate}
vanish and the profile is fixed. Its radius is
\(3v/[m(2m+v)]\sim3/(7m)\), attained at the third difference.

\begin{proposition}[Digit replication]\label{prop:replicate}
For any length, use its \(\lfloor n/2\rfloor\) differences
of opposite sorted digits. At any base and replication
factor \(s\ge2\), repeating every digit \(s\) times commutes with the sorted
multiset update exactly at height zero or when \(q_1+q_\theta=B\).
A cycle lifts phase by phase with its least period exactly
when every phase meets this condition. Every odd composite
\(N\ge9\) has at least \(\tau(N)\) nonconstant fixed points
for all sufficiently large bases in one progression, with
\(\tau(N)-1\) of them above a fixed positive radius cut,
where \(\tau(N)\) counts positive divisors.
The original mass of exactly replicated inputs from length
\(n<N\) is \(O_N(B^{n-N})\).
\end{proposition}

\begin{proof}
The provisional multiplicities in the borrowed list scale by
\(s\). The exact corrections change the number of digits
at least \(t\) by \(\one_{t=B-q_1}-\one_{t=q_\theta}\), for
\(1\le t<B\). Replication scales neither exceptional threshold
when it precedes subtraction, so the two orders differ by
\(1-s\) times this threshold defect. Threshold counts determine
the multiset. If the defect is nonzero, the resulting counts
at the two exceptional thresholds are not divisible by \(s\),
precluding a replicated successor. At height zero both
outputs are zero. Injectivity preserves the period of the
multiset cycle, and applying the numerical update to each
predecessor realizes the lifted words with that same period.

For each proper divisor \(1<n<N\), write \(m_n=(n+1)/2\) and
take \(B=m_n A\). The multiset
\[
 A-1,A,2A-1,2A,\ldots,(m_n-1)A-1,(m_n-1)A,B-1
\]
has \(q_i=(m_n-i)A\) and reproduces itself under the borrowed
update since the corrections cancel. The equal-gap fixed point at
\(B=m_n\) generates this family at \(B=m_nA\) by the paired covariance
of \cite[Theorem~4 and Section~6]{G4}.
Its \(s=N/n\) replicate has \((s-1)/2\) trailing zero
differences, giving distinct off-centre fixed points for
distinct divisors. Add the uniform-gap fixed point at \(N\)
and the nonconstant paired fixed profile from a small positive
multiple of \(v_\gamma-\overline v_\gamma\one\) at length \(N\).
The latter has only positive differences and is distinct from
all replicated families and from the centre. A common multiple
of the profile denominators and the \(m_n\) gives the progression.
There are at most \(O_N(B^n)\) ordered words in a replicated
locus, even allowing all permutations. The mass bound follows.
\end{proof}

In particular, at \(n=4^a-1\), \(a\ge2\), choose \(L\) even
in the paired sharpness example. At the same base the replicated
three-digit fixed point has radius exactly \((n-3)/[2(n+1)]\), attained
at the first difference, which tends to \(1/2\) as \(a\to\infty\).
Its directly replicated original mass is \(O_n(B^{3-n})\).

\section{Original pairing and chamber entrance}\label{sec:entrance}

\begin{theorem}[Initial pairing and reflected entrances]\label{th:sibling}
At odd length \(n\ge5\) the all-base limit of the original
positive-pairing probability is
\((5/8)^{\lfloor(n-1)/4\rfloor}\).
The initial chamber probabilities at 7, 9 and 11 digits
tend to \(5/16\), \(25/128\) and \(15/128\), respectively.
At each of the 47 strict lengths in \cite{G9}, the
probability that the initial gap is not positively paired, that the
gaps lie in \(\C\) from gap time one on, and that the orbit is captured
with \(Z_B>\delta_n\) by gap time \(T+1\) is
\(\asymp_n((T+1)/B)^r\) for \(T_0(n)\le T\le c_nB\).
Here \(\delta_n\) and \(c_n\) may be smaller than in
Theorem~\ref{th:clock}, and \(T_0(n)\) is a further constant.
\end{theorem}

\begin{proof}
Under the limiting gap density, the prefix sums
\(L_i=\sum_{j<i}g_j\) in normalized coordinates have the law
of the even order statistics of \(n\) auxiliary uniform
points, and \(R_i=1-L_{m-i}\) are the reflected suffix sums.
The continuum borrowed levels are \(1-L_i,1-R_i\).
With cross-list ties of probability zero, pairing holds
exactly when the two counting functions differ by at most
one at every threshold. A gap of two in the counts means
the \((i+1)\)-st arrival from one side preceded the \(i\)-th
from the other, precisely a level overlap.

Reveal the uniform points in increasing distance to the
nearer end of the unit interval, recording independent fair
left and right labels. Up to the midpoint the two counts
after \(a\) left and \(b\) right choices are
\(\lfloor a/2\rfloor,\lfloor b/2\rfloor\).
Reflection gives the counts beyond the midpoint as
\(m-1-\lfloor b/2\rfloor,m-1-\lfloor a/2\rfloor\).
After \(4j+1\) revealed points only \(a-b=\pm1\) is
allowed. From unit counts at those two states, the next
three additions give counts \((1,2,1)\), \((1,3,3,1)\), \((4,6,4)\)
in increasing order of allowed differences. The fourth
gives \((10,10)\). So the survival probability is multiplied
by \(10/16\) per block. The optional final two additions
have no deletion.
For the chamber also require \(L_i<R_i\) for
\(2\le i\le m-2\), using \(d_{m-i}=d_i\).
In the side-word recurrence put \(c_{0,0}=1\) and
\(c_{a,b}=c_{a-1,b}+c_{a,b-1}\) at valid positive-length
prefixes. Delete a prefix if the floor counts differ by more
than one, or if
\(\lfloor a/2\rfloor<i\le\lfloor b/2\rfloor\) for any such
\(i\). A failure beyond the midpoint corresponds to the
failure at \(m-i\) below it. The totals at 7, 9 and 11 are
40, 100 and 240 out of \(2^n\). Rounding and borrowing
have bounded integer errors, so the continuum rules give
the all-base limits.

At full positive height each interior difference \(v\) gives
the unordered borrowed pair \(\{v,B-1-v\}\).
Reflecting any subset of interior differences, resorting,
and keeping within the endpoints therefore gives identical
successor gaps and hence the same second numerical word.
Each reflected source has its own original multiplicity.

Use the interior starts constructed for the marked lower
bound in Section~\ref{sec:transport}. Their strict profile
can be chosen close to uniform, with all gaps, level
separations and \(d_2\) bounded below by positive multiples
of \(B\). Replace \(q_{k-1}\) by \(y=B-1-q_{k-1}\).
Since \(q_2-y=d_2+1\) and the original levels 2 and 3
are separated, \(y\) lies strictly between \(q_2\) and \(q_3\)
with margins of order \(B\). The new third difference is
\(y\), so reflection there recovers the original tuple.
The endpoints are preserved since these strict lengths
have \(k\ge4\). All gaps still have size comparable to \(B\).
Above a threshold between the original levels 2 and 3,
the four borrowed digits, aside from the singleton, now
consist of three from one column and one from the other.
This is incompatible with positive pairing. The common
successor gives permanent chamber residence from time one
and the same marked terminal fate. The injective reflection
and original weights give the lower bound. The upper
bound is the prefix statement in Theorem~\ref{th:clock}(1)
with \(j=1\).
\end{proof}

\begin{theorem}[Entrance peak and clock]\label{th:entrance}
Let \(S=\{0\}\cup\operatorname{orb}_P(1)\cup\operatorname{orb}_P(3)\),
and for an entrance gap \(g\) put
\[
 U(g)=\sum_{j\in S}g_j/B,\qquad
 D_P(g)=\max_{0\le t<r} R(P^t g).
\]
At a permanent chamber entrance,
\(|Z_B-D_P(g)|\le U(g)\).
At a strict length the gap time from that entrance to the
first cycle equals the transient of its relative root.
For an original start with \(g_0\in\C\), let \(\tau\le\infty\) be the
first time with \(g_\tau\notin\C\); when \(\tau<\infty\), \(g_\tau\) is
the image of the last chamber source. If \(|S|/m\le\delta/3\) and
\(0<\delta<1\), then
\[
 \limsup_{B\to\infty}
 \PP_{n,B}\Bigl\{g_0\in\C,\ \tau<\infty,\
 \max_{0\le t\le\tau}R(g_t)>\delta\Bigr\}
 \le (2kr+1)e^{-2n\delta^2/9}.
\]
This decays exponentially as \(n\to\infty\) with
\(r=o(n)\), including \(n=2^a\pm1\). The hypothesis gives
\(\delta\ge3|S|/(k+1)>2/k\), since \(k\ge3\), so by
Theorem~\ref{th:radius} the permanent terminal event
\(\{g_t\in\C\ \text{for all }t\ge0,\ Z_B>\delta\}\) is empty for
\(B>6/(\delta-2/k)\).
\end{theorem}

\begin{proof}
All corrections in \eqref{eq:Cupdate} are supported on
the invariant set \(S\). The physical and permuted gaps
agree off \(S\) while the chamber update applies and are
nonnegative with equal total on \(S\). Each suffix sum
can therefore change by at most \(BU(g)\) relative to the
permuted gaps. Over any \(r\) successive terminal times
the reference phases visit a full folded turn. This
proves the peak estimate, including each finite step
through the image of the last chamber source.

At a strict length a root recurring from time \(a\)
with period \(\ell\) on a permanent path must be loss-free
thereafter. Positive loss per period would eventually
make the complementary sum negative. With no losses,
the absolute root sum and its relative coordinates determine
its gaps, and the complement simply permutes. Thus the
full gap returns after \(\lcm(\ell,r)\) steps from time \(a\).
Any gap return implies a root return. If the entrance is the
first permanent one, no earlier gap lies on a chamber cycle.

Under the limiting original gap law, the sum of any fixed
\(j\) slots has distribution \(\mathrm{Beta}(2j,2m-2j)\)
for \(0<j<m\), equivalently the \(2j\)-th order statistic
of \(n\) uniforms. The binomial tail bounds
\cite[Theorem~2]{Hoeffding} give
\[
 \PP\{X-j/m>v\},\quad \PP\{X-j/m<-v\}\le e^{-2nv^2},
\]
since \((2j-1)/n<j/m<2j/n\).
The empty and full sums are deterministic. Apply this with
\(v=\delta/3\) to the \(kr\) rotated suffixes and to \(S\).
Except for the union of the tails, \(D_P\le\delta/3\)
and \(U\le2\delta/3\). The original gaps converge under
rounding and the relevant interior threshold boundaries
have limiting measure zero. This proves the asserted bounds.
\end{proof}

For a chamber terminal cycle let \(g_e\) be its first permanent
entrance, and let \(E\) be the event that \(g_e\) exists. For
\(\eta>0\) the comparison with \(D_P(g_e)\) gives the joint bounds
\[
 \begin{aligned}
 \PP_{n,B}(E,\ Z_B>\delta)
   &\le\PP_{n,B}(E,\ D_P(g_e)>\delta-\eta)+\epsilon_\eta,\\
 \PP_{n,B}(E,\ Z_B>\delta)
   &\ge\PP_{n,B}(E,\ D_P(g_e)>\delta+\eta)-\epsilon_\eta,
 \end{aligned}
\]
with \(\epsilon_\eta=\PP_{n,B}(E,\ U(g_e)>\eta)\). Conditional
probabilities given \(E\) would require dividing the error by
\(\PP_{n,B}(E)\).
Let \(t_e\) be the first permanent entrance time, so that \(g_e=g_{t_e}\). If \(q_k(g_t)>0\) for
\(0\le t<t_e\), then integer gap majorization \cite[Theorem~1]{G4} gives
\(U(g_e)\le L_s(g_0)+2t_e/B\), where \(L_s\) sums the
\(s=|S|\) largest normalized gaps. A full-height suffix of only
\(\ell\) sources gives \(U(g_e)\le L_s(g_{t_e-\ell})+2\ell/B\) instead.
Full height is needed already for the underlying rank-one estimate: at
\(n=9\), \(B=1000\) the source \(q=(675,671,336,0)\), with gaps
\((325,4,335,336,0)\), has successor gaps \((326,3,7,652,12)\), so the
largest gap rises from 336 to 652 in one step, violating largest-gap
majorization. (There \(S\) is the full slot set, so \(U=L_5=1\).)
A union over the
\(\binom ms\) subsets in the beta bound shows that
\(L_s(g_0)\) tends to zero in probability in the iterated
base-then-length limit for \(s=o(m)\).
In particular the entrance error is small for prefixes of
\(o(B)\) full-height steps in those limits. Estimating \(D_P\) at
arbitrary entrances requires its own access analysis.

\section{The integer root recurrence}\label{sec:root}

The autonomous map \eqref{eq:root} is defined on \(\Z^r\)
for each \(r\ge3\). If \(y_t=x_r(t)\) is the read at
time \(t\), put \(e_t=\sgn y_t\) and \(h_t=(-1-y_t)_+\).
Following that coordinate around a turn gives
\begin{equation}\label{eq:reads}
 y_{t+r}-y_t=-3e_t+2h_t+2h_{t+1}
       +\sum_{j=2}^{r-1}(h_{t+j}-e_{t+j}).
\end{equation}
The first \(r\) reads access the initial slots in order
\(r,r-1,\ldots,1\). Each read equals its initial coordinate
plus corrections determined by preceding reads. Subtracting
these corrections in order recovers the initial vector.
This is an integer bijection conjugating \(T_r\) to the
shift on consecutive reads with recurrence \eqref{eq:reads}.
In particular the least periods of the read word and state
orbit agree.

\begin{proposition}[Phase balance]\label{prop:balance}
On any loss-free root cycle of least period \(p\), the read
signs sum to zero in every class modulo \(\gcd(p,r)\).
\end{proposition}
\begin{proof}
Let \(d_*=\gcd(p,r)\), and let \(E_c\) be the sum of signs
at phases in class \(c\) modulo \(d_*\). Summing
\eqref{eq:reads} over that class, the left side cancels and
the loss-free right side gives
\[
 0=2E_c-E_{c+1}+(r/d_*)\sum_j E_j .
\]
The last term comes from a full turn of shifts. Summing in
\(c\) shows the total sign sum vanishes. Hence
\(E_{c+1}=2E_c\) cyclically, forcing every \(E_c=0\).
\end{proof}

\begin{theorem}[Physical guard certificate]\label{th:guards}
Fix a strict length. Suppose the nominal relative root
path at a chamber entry has transient \(a\) and loss-free
terminal period \(\ell\). Continue the nominal full gaps by
\eqref{eq:Cupdate}. The visit remains permanently in
\(\C\) if and only if the physical gap, height, pairing and
secondary guards pass at source phases
\(0,\ldots,a+\lcm(\ell,r)-1\).
For an integral invariant vector \(w\) constant on \(J\),
translate \((g,B)\) by \(s(w,\sum_j w_j)\).
The gap-sum identity is preserved by the translation, the nominal
reads are unchanged and the guard values
vary affinely with \(s\). All sufficiently large integer
\(s\) pass the finite test exactly when every slope is
nonnegative and every zero-slope guard already holds.
A positive near-uniform secondary-strict \(w\) can be
chosen with strictly positive slopes for all gap, height,
interlevel pairing and secondary inequalities.
\end{theorem}

\begin{proof}
After the root transient, all losses vanish. The full
nominal gaps then repeat after a common multiple of the
root and permutation periods, by conservation of the
root sum. The finite test checks exactly the sources
through the first such return, and hence every source
thereafter. If the physical visit stays in the chamber,
its update agrees with the nominal one, so the same
guards are necessary. The same repeated block gives
the actual peak.

The translated nominal trajectory is \(g_t+sw\) at every
time because the primary row, and hence the sign and loss,
are unchanged while \(Pw=w\). Write every tested inequality
as an affine integer form \(G\ge0\). Its translated value is
\(G(g_t,B)+sG_{\mathrm{hom}}(w,\sum_jw_j)\).
This proves the slope criterion and gives a computable
threshold for each positive slope. A rational strict
invariant profile can be mixed with uniform gaps and
cleared of denominators. Close enough to uniform, the
physical gaps, full height and interlevel comparisons
have positive slopes. Strictness gives positive slopes
for the secondary rows. The resulting physical periods
have a peak converging to the positive radius of
\(w/\sum_jw_j\).

Secondary strictness by itself can fail pairing.
At \(n=9\) and base 50, translate uniform gaps of size 10
by \(w=(1,1,1,10,1)\), with the base increased by \(14s\).
The vector is invariant, constant on \(J=\{0,1,2,4\}\)
and \(d_2(w)=9\). The nominal root stays zero. Along
this ray the second reflected level minus the third
direct level is \(B_s-2q_3-1=9-8s\), so pairing fails
for integers \(s\ge2\).
\end{proof}

\section{A graph for periods dividing \texorpdfstring{\(2r\)}{2r}}
\label{sec:graph}

\begin{theorem}[Periodic root graph]\label{th:graph}
Let \(G\) be the directed graph on the 21 states
\((u,e)\in\{-3,\ldots,3\}\times\{-1,0,1\}\) with an edge
\((u,e)\to(u-2e,f)\) whenever \(u-2e\in[-3,3]\) and \(f\le u\) if
\(e=1\), \(f\ge u\) if \(e=-1\), \(f=u\) if \(e=0\). For \(r\ge3\) call
\(x\in\Z^r\) \emph{covered} if it is a loss-free periodic point of
\(T_r\) with \(T_r^{2r}x=x\), that is, of least period dividing \(2r\).
Let \(y_t\) be the read of \(T_r^tx\), \(e_t=\sgn y_t\) and
\(u_t=\sum_{j=0}^{r-1}e_{t+j}\). The map
\(x\mapsto(u_t,e_t)_{0\le t\le r}\) is a bijection from covered points
onto the length-\(r\) walks in \(G\) that end at \((-u_0,-e_0)\). Its
inverse sets \(y_t=1+u_t-e_{t+1}\), \(-1\) or \(0\) according as
\(e_t=1\), \(-1\) or \(0\), for \(0\le t<r\), and recovers \(x\) from
these reads. Exactly 13 states of \(G\) occur on such walks across all
orders. All covered reads lie in \([-1,5]\). The number of covered
points, or read words with a specified initial phase, is
\[
 \begin{cases}
 2^r+1,&r\text{ odd},\\
 (2^{r/2}+1)^2,&r\text{ even}.
 \end{cases}
\]
A cycle of least period \(p\) contributes \(p\) points.
\end{theorem}

\begin{proof}
On a loss-free path put \(u_t=\sum_{j=0}^{r-1}e_{t+j}\).
Equation~\eqref{eq:reads} becomes
\begin{equation}\label{eq:lossfree}
 y_{t+r}-y_t=-2e_t+e_{t+1}-u_t .
\end{equation}
If \(y_{t+2r}=y_t\), Proposition~\ref{prop:balance} gives
\(e_{t+r}=-e_t\): for least period \(p\mid2r\), a class
modulo \(\gcd(p,r)\) has either one phase, with zero sign, or
two phases separated by \(r\).
Thus \(u_{t+r}=-u_t\) and \(u_{t+1}=u_t-2e_t\).
For \(e_t=1\), the opposite read is \(-1\), so
\(y_t=1+u_t-e_{t+1}\), requiring \(u_t\ge e_{t+1}\).
For \(e_t=-1\), \(y_t=-1\) and the opposite read is
\(1-u_t+e_{t+1}\), requiring \(u_t\le e_{t+1}\).
For \(e_t=0\) both reads are zero and \(u_t=e_{t+1}\).
These conditions also suffice for the paired reads.

At a positive maximum of \(u\), the previous step cannot
be a decrease. A stationary step has \(|u|\le1\).
An increase starts at \(u\le e_{t+1}\le1\), so reaches
at most 3. Antisymmetry gives \(|u|\le3\).
These are the states of \(G\), and its edges from \((u,e)\) go to
\((u-2e,f)\) in that set with
\begin{equation}\label{eq:edge}
 \begin{array}{c|ccc}
 e&1&-1&0\\ \hline
 \text{condition}&u\ge f&u\le f&u=f
 \end{array}
\end{equation}
The involution \(\iota(u,e)=(-u,-e)\) preserves the graph.
A covered root orbit determines a walk to its opposite state.
Conversely follow such a walk by its image under \(\iota\).
The \(r\) updates of the first coordinate give
\(u_0=\sum_{j=0}^{r-1}e_j\), and the same identity holds
starting at every phase. Assign each positive read its
value \(1+u_t-e_{t+1}\), each negative read \(-1\), and
each zero read 0. Equation~\eqref{eq:edge} ensures that
these values lie in \([-1,5]\) and satisfy
\eqref{eq:lossfree}. The read bijection gives the stated
inverse.

To count the walks, decompose the graph into components.
For odd \(u\) write
\(A=(1,1),\ B'=(1,0),\ C'=(1,-1),\ D=(3,1)\)
and use a bar for the image under \(\iota\).
The edges are
\[
 A\to\bar A,\overline{B'},\overline{C'},\quad
 B'\to A,\quad C'\to D,\quad D\to A,B',C',
\]
together with their images under \(\iota\).
Among the even states on closed walks, a nonzero central
state \((0,e)\) goes to its distinct outer state \((-2e,-e)\). An outer
state goes to either nonzero central state, and also has an edge to
\((0,0)\), which has only its loop. These are the 13 states on closed
walks. The eight remaining nominal states have no outgoing edges. In
the full graph, a nonzero central state also has an edge to the
dead-end state \((-2e,0)\). The odd component and the four-state
nonzero even component are strongly connected and
preserved by \(\iota\), so each of their states occurs
on some walk to its opposite of length at least 3.

For the adjacency matrix \(M\) with rows indexed by
sources, the count is \(\tr(M^r\iota)\), where \(\iota\) also
denotes the permutation matrix. On the odd component use
functions satisfying \(f(\iota(s))=\epsilon f(s)\), with
\(\epsilon=\pm1\). On the coordinates \(A,B',C',D\) the
matrix is
\[
 M_\epsilon=\begin{pmatrix}
 \epsilon&\epsilon&\epsilon&0\\
 1&0&0&0\\0&0&0&1\\1&1&1&0
 \end{pmatrix},
 \qquad
 \det(\lambda I-M_\epsilon)=\lambda^2
 (\lambda^2-\epsilon\lambda-\epsilon-1).
\]
The difference of the two power traces is \(2^r\).
In the nonzero even component each step switches classes.
For \(r=2q\), two steps give one walk from each state to
each of the two states in its own class. This contributes
\(4\cdot2^{q-1}=2^{q+1}\). For odd \(r\) it contributes
zero. The remaining state contributes one.
\end{proof}

The period condition is essential. For example, at \(r=8\)
the read word
\[
 (-1,-1,-1,5,-1,5,0,1,1,-1,1,-1,1,0)
\]
satisfies \eqref{eq:lossfree} modulo 14 and has least period
14, which does not divide 16.

\begin{proposition}[Six-step insertion]\label{prop:six}
For \(r=7\) and every \(r\ge9\), the map \(T_r\) has a loss-free cycle
of least period \(2r\), given by the construction in the proof. The
certified order-15 cycle of Lemma~\ref{lem:cert15} is one of them.
\end{proposition}

\begin{proof}
We construct, for each odd core order \(r_0\ge7\), each even core order
\(r_0\ge10\), and each \(j\ge0\), a cycle of least period \(2r\) at
\(r=r_0+6j\), by inserting \(j\) of the six-step blocks below in each
turn of the core; \(j=0\) already gives every order in the statement.
The certified order-15 cycle is the order-9 core with one six-step block
in each turn.
Represent the two turns by a pair of read words \(U,V\).
The order-7 and order-10 seeds, both starting at graph
state \((1,-1)\), are
\[
\begin{array}{r|l|l}
r_0&U&V\\ \hline
7&(-1,3,3,-1,1,1,-1)&(1,-1,-1,3,-1,-1,5)\\
10&(-1,3,3,-1,0,3,-1,-1,3,1)&(1,-1,-1,2,0,-1,1,1,-1,-1).
\end{array}
\]
The successive first-turn window sums are
\[
 (1,3,1,-1,1,-1,-3)\quad\text{and}\quad
 (1,3,1,-1,1,1,-1,1,3,1).
\]
Substitution in \eqref{eq:edge} verifies the words,
including the joins.
After three edges the first turn is at \((-1,-1)\).
There the two-edge loop with signs \((-,+)\) reads
\((-1,3)\), and its opposite loop reads \((3,-1)\).
Repeating these loops gives the cores of each indicated
order. There is also a six-edge loop at \((1,-1)\)
with successive sums \((1,3,1,-1,1,-1)\) and signs
\((-,+,+,-,+,-)\). It and its opposite read
\[
 \begin{aligned}
 b_U&=(-1,3,3,-1,3,-1), & b_V&=(1,-1,-1,3,-1,1).
 \end{aligned}
\]
Prefixing these loops preserves the endpoints of the two
walks and therefore their joins.
The odd seed has exactly one 5 and the even seed exactly
one 2. These markers stay unique under both insertions,
so the word of length \(2r\) has no smaller period.
At \(r_0=9,j=1\) the first \(r\) reads recover the phase
\((2,2,1,0,3,0,3,2,-1,0,3,0,3,2,-1)\) of the order-15 record.
\end{proof}

The restriction to \(r=7\) and \(r\ge9\) belongs to the construction,
not to the graph. Enumerating the graph walks shows that such roots also
exist at orders 3--6 and 8; at \(r=8\), for example, there are 18 cycles
of least period 16 (Section~\ref{sec:verify}).

\begin{proposition}[Two possible final entries]\label{prop:doors}
A positive-time first entry into a cycle of
Theorem~\ref{th:graph} is loss-free on its last step.
The two possible unordered pairs of cyclic and entering
reads at the same target are \(\{0,3\}\) and \(\{-1,5\}\).
Each phase whose predecessor on the cycle has one of
these reads has exactly one off-cycle one-step predecessor.
\end{proposition}
\begin{proof}
Given the last coordinate \(y\) of a predecessor, its
other slots are recovered uniquely from target slots
\(2,\ldots,r\) using \eqref{eq:root}. The first target
slot is \(1-y\) for \(y\le-2\), 2 for \(y=-1\), 0 for
\(y=0\), and \(y-3\) for \(y\ge1\). Values written to the first slot
from covered reads lie in \([-2,2]\) and have only loss-free preimages.
The only multiple inverse values are 0 and 2 and their
two reads are the claimed pairs. Each phase on a cycle
has exactly one predecessor on that cycle.
\end{proof}

\section{Off-centre basin probability at order 15}
\label{sec:order15}

\begin{theorem}[Order-15 survival]\label{th:order15}
For each of the twelve lengths
\[
 \begin{gathered}
 99,151,217,331,993,1057,2979,3641,4681,10923,32767,32769 ,
 \end{gathered}
\]
there is \(\delta_n>0\) with
\(\liminf_{B\to\infty}\PP_{n,B}(Z_B>\delta_n)>0\),
through all integer bases. These are all lengths of
folded order 15 and all are strict.
\end{theorem}

\begin{lemma}[Order-15 reference records]\label{lem:cert15}
The four integer records supplied with the paper satisfy
the insertion conditions in Appendix~\ref{app:retirement}.
Each reaches the same loss-free period-30 root with first phase
\[
 (2,2,1,0,3,0,3,2,-1,0,3,0,3,2,-1).
\]
The fourth record has determinant of absolute value
\[1957661336689015451135390789227234983936000.\]
For every integer \(L\ge795308\) it yields at least
\(795308^{-15}L^{15}\) distinct roots in \([-L,L]^{15}\).
Their lifetime norm \(\sup_t\|x(t)\|_\infty\) is at most \(L\), their
displacement from the start is at most \(2L\) in maximum norm, and
their total loss is at most \((1933060/397654)L\).
\end{lemma}

\begin{proof}[Proof of Theorem~\ref{th:order15}]
The candidate lengths divide \(2^{15}-1=7\cdot31\cdot151\)
or \(2^{15}+1=3^2\cdot11\cdot331\).
The divisors at least 7 of lower folded order are
\(7,9,11,31,33\). For eleven of the remaining lengths,
take integer offsets zero on \(J\) and one elsewhere.
At 99, take zero on \(J\) and offsets
\((1,1,1,0,2,3)\) on the orbits with least positive
slots \((3,5,9,11,15,33)\), respectively.
In each case the offsets \(w_j\) are invariant and
all secondary rows have minimum slack one. This is
the finite check in Section~\ref{sec:verify}.
Let \(b=1+\sum_jw_j\), \(Q=mb+\sum_jw_j\), and
\(\pi_j=(b+w_j)/Q\). This normalized profile is
secondary-strict with \((m+1)\min_j\pi_j-1=1/Q\).
The latter inequality gives a positive homogeneous interlevel margin at
\(\pi\), which persists under sufficiently small normalized
perturbations. In fact for
the unborrowed profile the largest entry of level
\(i+1\) is at most \(1-(i+1)\min_j\pi_j\), while
the smallest of level \(i\) is at least
\((m-i)\min_j\pi_j\). The normalized radius of the
profile is positive.

By Lemma~\ref{lem:cert15} the hypotheses of strict source transport
\cite[Proposition~7]{G10} hold at \(r=15\): for every \(L\ge795308\)
there are at least \(795308^{-15}L^{15}\) distinct roots in
\([-L,L]^{15}\) whose lifetime norm and total loss are at most
\((1933060/397654)L\). That proposition gives the conclusion at every
strict length of folded order 15, through all sufficiently large integer
bases, with no congruence restriction on \(B\). Its proof takes
\(L=\lfloor B/D_n\rfloor\) for a sufficiently large constant \(D_n\),
lets the anchor and all but one complementary slot vary over
\(\gg_nL\) integers around \(B\pi\), for a rational strict profile
\(\pi\) with a positive pairing margin such as the one above, and uses
the root norm and loss bounds to keep every nominal gap within
\(O_n(L+1)\) of \(B\pi\). All guards then hold, and the resulting
cycles have peak bounded away from zero. These cycles lie in \(\C\), so
by Theorem~\ref{th:radius} the cut obtained in this way satisfies
\(\delta_n\le2/k\); at \(n=32767\) this is about \(1.2\cdot10^{-4}\).
\end{proof}

\section{Scope and open questions}\label{sec:open}

The growing-time exponent in Theorem~\ref{th:clock}
is a chamber-capture exponent. For arbitrary prefixes
Theorem~\ref{th:clock} gives the fixed-word-time bound (4) and, in (1),
a fixed number of unconstrained initial steps followed by chamber
residence; the one-step entrance theorem covers a specified permanent
chamber event. The
peak estimate gives an observable at a permanent entrance,
whose original distribution still has to be determined.
The paired radius bound allows all secondary comparison
failures. Its large-length consequences use iterated
limits, with the base tending to infinity separately at
each length. The zero-cone bound locates a failure of
pairing or a deep breach near an energy maximum. It gives
no estimate of the original probability of reaching that
window. Replication and the periodic root graph give the
existence of the corresponding cycles, with no weights. The positive
basin construction uses full rank, linear lifetime bounds and exact
guard margins. The order-15 result extends the positive-length list of
\cite{G9}. The unrestricted distribution of terminal periods and the
size of the off-centre share at general fixed length remain open.

\subsection*{Macroscopic mass and strictness}

The strictness and survival conjecture of \cite[Section~7]{G10}
asks that, at strict odd lengths \(n\ge7\), there be a cut
\(\delta>0\) with \(\liminf_B\PP_{n,B}(Z_B>\delta)>0\) over
all bases, and that at non-strict lengths the probability
tend to zero at every fixed positive cut.
If the class limits of the peak distribution exist, set
\[
 d_\beta=\lim_{a\downarrow0}\ 
  \lim_{\substack{B\to\infty\\B\equiv\beta\pmod m}}
  \PP_{n,B}(Z_B>a)
\]
through continuity cuts \(a>0\). In that formulation the
conjecture asks for \(\min_\beta d_\beta>0\) at strict
lengths and \(\max_\beta d_\beta=0\) at non-strict lengths.
Theorem~\ref{th:order15} establishes the original all-base
lower statement at order 15 without requiring class limits.

\subsection*{Terminal labels and the peak law}

At fixed odd length \(n\ge3\), put
\(v_*=(k,k-1,\ldots,1)\) and \(X=mq-Bv_*\).
For \(\beta=B\bmod m\), bounded offsets lie in
\(\Lambda_\beta=\{X\in\Z^k:X\equiv-\beta v_*\pmod m\}\).
Write the exact paired covariance of \cite[Theorem~4 and Section~6]{G4}
in base-free form as \(\Phi_k(X)=mF_B((X+Bv_*)/m)-Bv_*\), where \(F_B\)
is the difference update and \(B\equiv\beta\pmod m\) is sufficiently
large. The shift \((B,q)\mapsto(B+m,q+v_*)\) of \cite{G4} preserves both
\(X\) and the successor offset. For each fixed offset, the leading
interlevel margins imply positive pairing at sufficiently large base. A
cycle of fixed offset vectors is therefore a physical cycle throughout a
sufficiently large part of its class, and every family with the same
fixed physical offset cycle gives that label. The gcd conclusion of
\cite[Theorem~4]{G4} assumes a fixed Jordan-measurable conditioning
region of positive limiting weight in the strict simplex, all but
vanishing conditional mass with entirely paired difference orbits, and
proper joint limits along every class of a finite modulus \(M\); it then
asserts that the limits depend only on the base modulo \(\gcd(M,m)\).

It is open at general fixed length whether the original
probability of each exact bounded-offset cycle label
converges in each base class, and whether the physical
all-phase peak distribution converges, with a chamber or
exterior mark where the chamber is defined. The limiting
individual cycle weights, if they exist, form a
subprobability. Identifying its missing mass with
\(d_\beta\) requires the additional exclusion of
intermediate radii. One must control the band
\(K_0<BZ_B\), \(Z_B\le a\) by taking the class base
limit first, then \(a\downarrow0\) through continuity
cuts, and \(K_0\to\infty\). Finiteness of the
positive-weight support is a separate question.
A joint arrival law also requires the original
position, arrival time and entry phase to be retained.

\subsection*{Transfer in the folded order}

The six-step family constructs least-period-\(2r\) roots explicitly at
order 7 and every order from 9. Enumerating the graph walks shows that
such roots also exist at orders 3--6 and 8 (Section~\ref{sec:verify}).
To obtain positive basin mass at the strict lengths of a given folded
order \(r\) by \cite[Proposition~7]{G10}, it suffices to exhibit at least \(c_rL^r\) distinct roots in
\([-L,L]^r\) for every sufficiently large \(L\), with lifetime norm and
total loss \(O_r(L)\). Their terminal cycles may differ. A transfer construction valid across
unbounded orders is open. An immediate predecessor or
a single guarded trajectory supplies a smaller set of
sources than this estimate.

\section{Computations and verification}\label{sec:verify}

From a fresh copy of the ancillary directory run
\texttt{bash reproduce.sh}. The command verifies SHA-256
hashes of the inputs and compares the output of four
standard-library Python replays to their recorded exact
outputs. The replay was checked with CPython 3.13.15 and uses exact integer
and rational arithmetic. The computational inputs to the theorems are
as follows.

\begin{enumerate}
\item \emph{Marked capture at orders 3--13.}
\texttt{recon14\_physical/verify.py} independently reads the
strict profiles from \cite{G9} and three archived JSONL
datasets. It checks their 47 profile margins, exact
orders and peaks, and the full-rank source records of
Lemma~\ref{lem:p9}. The reference trajectories, terminal
periods, ray signs, rank, and time, spread and loss
coefficients are recomputed from integer updates.
The same data enter Theorem~\ref{th:sibling}.
The replay also gives finite literal controls for
the paired update, bounded inverses, the 25-digit ray
seeds and replication. On the two complementary orbit offsets
parametrizing \(W\) at \(n=229,283,1069\), the replay finds three
comparison rows spanning dimension two with a strictly positive linear
combination equal to zero. Nonnegative slacks therefore force both
offsets to vanish, and orbit averaging gives \(K=\{0\}\).

\item \emph{Order 15.}
\texttt{recon\_root15/replay.py} checks the two frozen
JSONL inputs and reconstructs the four reference
paths and insertion records of Lemma~\ref{lem:cert15}
in 11,240,158 updates. It uses the slot update and
checks the resulting reads against \eqref{eq:reads}.
It verifies each first nonpositive read, the
predecessor intervals and their return on retired
coordinates, the triangular determinants, the common
first period-30 entry and the lifetime budgets.
It enumerates the folded-order-15 lengths and checks
all comparison slacks and the positive peak of the
normalized witnesses in Theorem~\ref{th:order15}.

\item \emph{Independent finite controls.}
\texttt{recon15/check.py} checks the radius allowances,
counts of valid side words, reflected entrances, the
small-time geometry, the recurrence and the guard
translation against literal or exact examples.
\texttt{recon\_carriers/verify.py} checks the graph counts
at orders 3 through 80, its covered read words and one-step
predecessors through order 13, and ternary sign words
through order 9. (The replay output calls these words marked words.) It records the
least periods of all covered read words
through order 13: cycles of least period \(2r\) occur at every order
from 3 to 13, for example 18 of them at order 8. It also checks the
six-step construction
and the example of period 14 at order 8. The general
rank, probability, graph and guard results follow from
the written arguments. The arrival laws themselves
are the results of \cite{G3,G5}.
\end{enumerate}

\appendix
\section{From a finite reference path to many starts}
\label{app:retirement}

Here is the insertion argument underlying the two full-rank
replays. It is the retirement construction of
\cite[Section~4, Lemma~7, and Appendix~B]{G9}.
Let \(S_r\) rotate the last slot to the first. Label the
coordinates by their original slots and follow those
labels through successive rotations.
In a reference path \(x(t)\), each insertion interval
\([s_i,t_i)\) has length divisible by \(r\) and lies before
the terminal cycle. The intervals are ordered and disjoint.
A label is called retired after its first nonpositive read.
The ray
\[
 u_i=S_r^{-s_i}(x(s_i)-x(t_i))
\]
has nonnegative entries, vanishes on labels read
nonpositively before \(t_i\), and has positive entry on
the label that retires next. In retirement row order
the columns \(u_i\) thus form a lower triangular matrix
with positive diagonal. The complete reference cycle
return is loss-free. The finite records verify these
properties directly.

Take the start \(x(0)+\sum_i a_i u_i\) with
\(0\le a_i<M\) integral. In any prefix before the
interval for a still unconsumed ray, that ray is zero
on every label whose read is nonpositive. At positive
reads, adding remaining copies keeps the read positive.
Therefore the sign and loss stay those of the reference
and the perturbation is transported by \(S_r\).
During interval \(i\), one traversal from the
state perturbed by \(S_r^{s_i}p\) ends at
\(x(s_i)+S_r^{s_i}(p-u_i)\).
Insert \(a_i\) such traversals before proceeding along
the reference interval. At each inserted traversal
the same sign argument applies to the remaining rays.
Induction removes all added copies and reaches the
reference terminal cycle. The determinant gives \(M^r\)
distinct starts.

Reference time through a complete terminal return,
plus the sum of the interval lengths, bounds the
time coefficient after multiplication by \(M\).
The reference total loss plus the sum of the interval
losses likewise bounds the loss coefficient. The
reference lifetime norm plus the sum of ray norms bounds
the norm coefficient. The spread bound uses the same
argument with the range including zero, which is
subadditive. This proves the family in Lemma~\ref{lem:p9}
by taking \(M=L\).

For the fourth order-15 record the reference lifetime
maximum norm is 256429 and the sum of ray norms is
141225. The reference loss is 1256253 and the intervals
have total loss 676807. The time through one terminal
return is 2428969, and the intervals have total length
1308120. The resulting norm, loss and time coefficients
are 397654, 1933060 and 3737089, respectively. Take
\(M=\lfloor L/397654\rfloor\). For \(L\ge795308\),
\(M\ge L/795308\). The lifetime norm is at most \(L\)
and the displacement from each start is at most \(2L\).
This gives the bounds of Lemma~\ref{lem:cert15}.

\section{Small-time tails of the arrival laws}\label{app:tails}

Write \(F_5\) for the limiting distribution function of
normalized first cyclic-word time at five digits and
\(F_{7,\beta}\) for the corresponding seven-digit
distribution function along base parity \(\beta\).
These are the original-word laws of \cite{G3,G5}.

\begin{proposition}[Small-time tails]\label{prop:tails}
As \(t\downarrow0\),
\[
 F_5(t)=10t^2+O(t^3),\qquad
 F_{7,\beta}(t)\ge c t^3\log(1/t)\quad(\beta=0,1)
\]
for a common \(c>0\) in the second bound. In particular
\(F_{7,\beta}(t)/t^3\) is unbounded at zero for each parity.
\end{proposition}

\begin{proof}[Five digits]
In original differences \((x,y)\), \(0<y<x<1\), the
density is \(120(1-x)y(x-y)\). The initial fold
\cite[Section~1.4]{G3} has four sheets with values of
\((U,V)\) equal to
\[
 (1-y,1-x-y),\quad (1-y,x-y),\quad
 (x,x-y),\quad (x,x+y-1)
\]
in the regions \(x<1/2\), \(x\ge1/2,\ x+y<1\),
\(x+y\ge1,\ y<1/2\), and \(y\ge1/2\), respectively.
The shear gives \(w=V-j(1-U)\), with
\(j=\max(0,\lceil(V-1/2)/(1-U)\rceil)\).
The clock is \(2/3\) times the range of \(0,A,C'\),
where \(A=U+w-1\) and \(C'=2U-w-1\).
For \(h=3t/2\), its sublevel is the hexagon
\(|A|,|C'|,|A-C'|\le h\).
It drives \((U,w)\) to \((2/3,1/3)\), so \(U<3/4\)
for small \(t\). Each sheet satisfies \(V\le2U-1\),
which gives \(j=0\) there.

The four inverse fold centres are
\[
 (1/3,1/3),\quad (2/3,1/3),\quad (2/3,1/3),\quad (2/3,2/3).
\]
The fold Jacobians have absolute value one, and the
\((U,w)\)-to-\((A,C')\) Jacobian has absolute value three.
At the interior centre the middle two sheets each take
the \(C'\ge0\) half of the hexagon: on them
\(C'=1-x-y\) and \(C'=x+y-1\), respectively.
The full hexagon has area \(3h^2\).
The two half-hexagons therefore contribute total
original area \(h^2\) at density \(40/9+O(t)\).
On the two outer sheets, \(x-y=O(t)\) and the area
is \(O(t^2)\), so their mass is \(O(t^3)\).
This gives \((40/9)h^2=10t^2\) as the leading term.
\end{proof}

\begin{proof}[Seven digits]
In the ordered absolute radii \((a,b,c)\) the leading
map of \cite[Sections~1 and~3]{G5} is
\[
 H(a,b,c)=\ord(a,|a+b-1|,|b+c-1|).
\]
Use \(v=a-b,\ w=a+b+c-1\). On the branch
\(0<c<b<a<1/2,\ 2a+b>1,\ w>0\) the largest radius
is preserved and the other coordinates update by
\((v,w)\mapsto(w,v-w)\).
Indeed \(a+b,b+c<1\) and the second largest
output is \(1-b-c<a\).
The inverse on \((v,w)\) is \(M(v,w)=(v+w,v)\)
of determinant \(-1\).

Take the box \(1/2-2\epsilon<a<1/2-\epsilon\),
\(\epsilon<v,w<5\epsilon/4\), and its \(j\)-fold
inverse boxes for \(0\le j\le\lfloor\log_2(1/(128\epsilon))\rfloor\),
with \(0<\epsilon\le1/1024\). All of the inverse steps
have \(0<v,w\le5/512\). With
\(b=a-v,\ c=1-2a+v+w\), these bounds keep the paths
on the stated branch. Each box has volume
\(\epsilon^3/16\). Their \(v\)-intervals are contained in
\((F_{j+2}\epsilon,5F_{j+2}\epsilon/4)\), where
\(F_1=F_2=1\) are the Fibonacci numbers.
Successive ratios are at least \(3/2\), so the boxes
are pairwise disjoint.

Map back to original differences by
\(\xi=((1+a)/2,(1+c)/2,(1-b)/2)\),
an affine inverse of the initial ordered-radius fold
with absolute determinant \(1/8\) in these coordinates.
All four original normalized gaps exceed \(1/8\).
Thus each inverse box contributes mass \(\gg\epsilon^3\).
Through the target step these orbits have largest
radius below \(1/2\), whereas the entrance core in
\cite[Theorem~1]{G5} requires \(a\ge1/2\).
The original \(\xi\) is also outside that core because
its first difference is below \(3/4\).
At the target the corresponding differences
\(\tfrac12(1+a,a+b,b+c)\) lie within \(8\epsilon\)
of \((3/4,1/2,1/4)\), so the gaps lie within
\(16\epsilon\) of \(1/4\).
The continuous majorization of \cite[Theorem~1]{G4}
keeps all later gaps within this interval.
By \cite[Proposition~3 and Theorem~1]{G5}, almost every source enters
the core and reaches the canonical phase, whose differences thus
lie within \(48\epsilon\) of centre.
Both clock functions in either row of that law vanish
at centre and have gradient of \(\ell^1\)-norm at most
two. Both parities therefore place the full box mass
in \([0,96\epsilon]\).
Summing the disjoint boxes and taking \(\epsilon=t/96\)
gives the lower bound.
\end{proof}

\section*{Tools}
The research, computations and drafting were carried out with substantial assistance from AI systems (OpenAI Codex and Anthropic Claude), under the author's direction and close supervision.
The exact replays use Python, and the paper was typeset in AMS-\LaTeX\ with Tectonic.
The author takes responsibility for the contents.

\begin{thebibliography}{RRJS2}
\raggedright
\bibitem[AP]{AP}
S. Akiyama and A. Peth\H{o},
\emph{Discretized rotation has infinitely many periodic orbits},
Nonlinearity \textbf{26} (2013), 871--880.
\href{https://doi.org/10.1088/0951-7715/26/3/871}{doi:10.1088/0951-7715/26/3/871};
\href{https://arxiv.org/abs/1206.3868}{arXiv:1206.3868}.

\bibitem[COST]{COST}
E. Chen, K. Ono, R. E. Schwartz and D. S. Thakur,
\emph{Four-digit Kaprekar dynamics in odd bases},
J. Integer Seq. \textbf{29} (2026), Article 26.4.7.
\href{https://cs.uwaterloo.ca/journals/JIS/VOL29/Ono/ono3.pdf}{Article 26.4.7}.

\bibitem[Dahl]{Dahl}
C. D. Dahl, \emph{Coarse-grained drift fields and attractor-basin entropy in Kaprekar's routine},
Entropy \textbf{28} (2026), Article 92.
\href{https://doi.org/10.3390/e28010092}{doi:10.3390/e28010092}.

\bibitem[DD]{DD}
I. Daubechies and R. DeVore,
\emph{Approximating a bandlimited function using very coarsely quantized data:
a family of stable sigma-delta modulators of arbitrary order},
Ann. of Math. (2) \textbf{158} (2003), 679--710.
\href{https://doi.org/10.4007/annals.2003.158.679}{doi:10.4007/annals.2003.158.679}.

\bibitem[DZ]{DZ}
P. Devlin and T. Zeng, \emph{Maximum distances in the four-digit Kaprekar process},
Integers \textbf{21} (2021), Article A97.
\href{https://math.colgate.edu/~integers/v97/v97.pdf}{Article A97}.

\bibitem[G1]{G1}
R. Galyavin, \emph{Cycles of the four-digit Kaprekar map in even bases}, 2026.
\href{https://doi.org/10.5281/zenodo.24358739}{doi:10.5281/zenodo.24358739}.

\bibitem[G2]{G2}
R. Galyavin, \emph{Arrival times for the four-digit Kaprekar map in even bases}, 2026.
\href{https://doi.org/10.5281/zenodo.24358744}{doi:10.5281/zenodo.24358744}.

\bibitem[G3]{G3}
R. Galyavin, \emph{A common arrival-time law for five-digit Kaprekar dynamics}, 2026.
\href{https://doi.org/10.5281/zenodo.24358746}{doi:10.5281/zenodo.24358746}.

\bibitem[G4]{G4}
R. Galyavin, \emph{Majorization in odd-digit Kaprekar dynamics and the nine-digit arrival law}, 2026.
\href{https://doi.org/10.5281/zenodo.24358748}{doi:10.5281/zenodo.24358748}.

\bibitem[G5]{G5}
R. Galyavin, \emph{Seven-digit Kaprekar dynamics: global law, two clocks, parity}, 2026.
\href{https://doi.org/10.5281/zenodo.24358752}{doi:10.5281/zenodo.24358752}.

\bibitem[G6]{G6}
R. Galyavin, \emph{Linear-time access for eleven-digit Kaprekar dynamics}, 2026.
\href{https://doi.org/10.5281/zenodo.24358756}{doi:10.5281/zenodo.24358756}.

\bibitem[G7]{G7}
R. Galyavin, \emph{Arithmetic selection of arrival times in eleven-digit Kaprekar dynamics}, 2026.
\href{https://doi.org/10.5281/zenodo.24371745}{doi:10.5281/zenodo.24371745}.

\bibitem[G8]{G8}
R. Galyavin, \emph{Eleven-digit Kaprekar arrival laws modulo twelve}, 2026.
\href{https://doi.org/10.5281/zenodo.24396437}{doi:10.5281/zenodo.24396437}.

\bibitem[G9]{G9}
R. Galyavin, \emph{Off-centre Kaprekar cycles with positive basin density}, 2026.
\href{https://doi.org/10.5281/zenodo.24399342}{doi:10.5281/zenodo.24399342}.

\bibitem[G10]{G10}
R. Galyavin, \emph{Macroscopic cycles in odd-length Kaprekar dynamics: a chamber criterion and a residue-transfer rule}, 2026.
\href{https://doi.org/10.5281/zenodo.24408723}{doi:10.5281/zenodo.24408723}.

\bibitem[G11]{G11}
R. Galyavin, \emph{The arithmetic spectrum of Kaprekar lengths}, 2026.
\href{https://doi.org/10.5281/zenodo.24419878}{doi:10.5281/zenodo.24419878}.

\bibitem[H]{Hoeffding}
W. Hoeffding, \emph{Probability inequalities for sums of bounded random variables},
J. Amer. Statist. Assoc. \textbf{58} (1963), 13--30.
\href{https://doi.org/10.1080/01621459.1963.10500830}{doi:10.1080/01621459.1963.10500830}.

\bibitem[Hof]{Hoffman52}
A. J. Hoffman, \emph{On approximate solutions of systems of linear
inequalities}, J. Research Nat. Bur. Standards \textbf{49} (1952),
263--265.
\href{https://nvlpubs.nist.gov/nistpubs/jres/049/4/V49.N04.A05.pdf}{NIST copy}.

\bibitem[KD1]{KDO}
A. Kay and K. Downes-Ward, \emph{Fixed points and cycles of the Kaprekar
transformation: 1. Odd bases},
J. Integer Seq. \textbf{25} (2022), Article 22.6.7.
\href{https://cs.uwaterloo.ca/journals/JIS/VOL25/Kay/kay5.html}{Article 22.6.7}.

\bibitem[KD2]{KDE}
A. Kay and K. Downes-Ward, \emph{Fixed points and cycles of the Kaprekar
transformation: 2. Even bases}, 2024.
\href{https://arxiv.org/abs/2408.12257}{arXiv:2408.12257}.

\bibitem[M]{Massart}
P. Massart, \emph{The tight constant in the Dvoretzky--Kiefer--Wolfowitz inequality},
Ann. Probab. \textbf{18} (1990), 1269--1283.
\href{https://doi.org/10.1214/aop/1176990746}{doi:10.1214/aop/1176990746}.

\bibitem[PGR]{PGR}
E. Perez Gonzalez and J. D. Reiss,
\emph{Idle tone behavior in sigma delta modulation},
Audio Engineering Society 122nd Convention, Vienna, 2007,
Convention Paper 7108.
\href{https://joshreiss.github.io/documents/2007/PerezReiss-AES122.pdf}{Author's copy}.

\bibitem[RRJS]{RRJS}
D. Reefman, J. Reiss, E. Janssen and M. Sandler,
\emph{Description of limit cycles in sigma--delta modulators},
IEEE Trans. Circuits Syst. I Regul. Pap. \textbf{52} (2005), 1211--1223.
\href{https://doi.org/10.1109/TCSI.2005.849134}{doi:10.1109/TCSI.2005.849134}.

\bibitem[RRJS2]{RRJS2}
D. Reefman, J. D. Reiss, E. Janssen and M. Sandler,
\emph{Description of limit cycles in feedback sigma delta modulators},
Audio Engineering Society 117th Convention, San Francisco, 2004,
Convention Paper 6280.
\href{https://joshreiss.github.io/documents/2004/ReefmanReiss-AES117-2004.pdf}{Author's copy}.

\bibitem[T]{Thakur}
D. S. Thakur, \emph{Kaprekar phenomena},
Proc. Ropar Conf., RMS Lecture Notes Ser. \textbf{26} (2019), 1--10.
\href{https://people.math.rochester.edu/faculty/dthakur2/kaprekar.pdf}{Author's copy}.

\bibitem[W]{Walden}
B. L. Walden, \emph{Searching for Kaprekar's constants: algorithms and
results}, Int. J. Math. Math. Sci. \textbf{2005} (2005), no.~18,
2999--3004.
\href{https://doi.org/10.1155/IJMMS.2005.2999}{doi:10.1155/IJMMS.2005.2999}.
\end{thebibliography}
\end{document}
